\input{../tex-inputs/latex-leseansicht-vorspann}

\section[Rainer Maria Rilke: Widmungsexemplar Lieder der Mädchen für Arthur Schnitzler, 14. 7. 1899]{L04544 Rainer Maria Rilke: Widmungsexemplar Lieder der Mädchen für Arthur Schnitzler, 14. 7. 1899}
\nopagebreak\mylabel{L04544v}
\rehead{ }\normalsize\beginnumbering\briefempfaengerindex{Schnitzler, Arthur@\textsc{Schnitzler, Arthur}!zzzRilke, Rainer Maria@\emph{von Rainer Maria Rilke}!1899-07-141@{14. 7. 1899}|(be}
\toendnotes[C]{\smallbreak\pagebreak[2]}
\correspDesc{Versand  durch Rainer Maria Rilke am 14. 7. 1899 in Schmargendorf
\newline{}Erhalt  durch Arthur Schnitzler im Zeitraum [15. 7. 1899 – 19. 7. 1899?] in Wien}\toendnotes[C]{\smallbreak}
\buchAlsQuelle{\emph{Rainer Maria Rilke und Arthur Schnitzler. Ihr Briefwechsel.}. Mit Anmerkungen herausgegeben von Heinrich Schnitzler In: \emph{Wort und Wahrheit}, Jg. 13, Nr. 4, April 1958, S. 282–298, hier S. 295.}\toendnotes[C]{\smallbreak}
\pstart
           \noindent{}{\pb}\pwindex{Rilke, Rainer Maria 4.\,12.\,1875 Prag – 29.\,12.\,1926 Montreux@\textsc{Rilke, Rainer Maria} (4.\,12.\,1875 Prag – 29.\,12.\,1926 Montreux)!Lieder der Mädchen@\strich\emph{Lieder der Mädchen}|pw}\label{K_L04539-1v}\edtext{Herrn Dr. Schnitzler}{\lemma{\textnormal{\emph{Herrn Dr. Schnitzler}}}\Cendnote{\textnormal{Die Widmung stand in einem Separatdruck des Texts aus der
                  Zeitschrift \emph{Pan}\pwindex{Pan@\emph{Pan}|pwk} (Rainer Maria Rilke\pwindex{Rilke, Rainer Maria 4.\,12.\,1875 Prag – 29.\,12.\,1926 Montreux@\textsc{Rilke, Rainer Maria} (4.\,12.\,1875 Prag – 29.\,12.\,1926 Montreux)|pwk}: \emph{Lieder der Mädchen}\pwindex{Rilke, Rainer Maria 4.\,12.\,1875 Prag – 29.\,12.\,1926 Montreux@\textsc{Rilke, Rainer Maria} (4.\,12.\,1875 Prag – 29.\,12.\,1926 Montreux)!Lieder der Mädchen@\strich\emph{Lieder der Mädchen}|pwk}. In: \emph{Pan}\pwindex{Pan@\emph{Pan}|pwk}, Jg. 4, H. 4, {[}April{]} 1899, S. 209–212.) Dieser befand sich 1958 im Besitz
                  von Heinrich Schnitzler\pwindex{Schnitzler, Heinrich 9.\,8.\,1902 Hinterbrühl – 12.\,7.\,1982 Wien@\textsc{Schnitzler, Heinrich} (9.\,8.\,1902 Hinterbrühl – 12.\,7.\,1982 Wien)|pwk}. Der Verbleib ist ungeklärt.}}}\label{K_L04539-1} mit vielen verehrungsvollen Grüßen\pend
           \pstart \spacefill\mbox{Rainer Maria Rilke.}\pend{}
\pstart
           Berlin-Schmargendorf\oindex{Schmargendorf@\textbf{Schmargendorf}|pw}, am 14. July 99.\pend
           \selectlanguage{ngerman}\endnumbering\briefempfaengerindex{Schnitzler, Arthur@\textsc{Schnitzler, Arthur}!zzzRilke, Rainer Maria@\emph{von Rainer Maria Rilke}!1899-07-141@{14. 7. 1899}|)be}\mylabel{L04544h}
\begin{anhang}
\end{anhang}\newcommand{\dateiname}{L04544}\newcommand{\titel}{Rainer Maria Rilke: Widmungsexemplar Lieder der Mädchen für Arthur Schnitzler, 14. 7. 1899}\newcommand{\editorInnen}{Herausgegeben von Jahnke, SelmaMüller, Martin Anton}\input{../tex-inputs/latex-leseansicht-abspann}
